% NAF 5161, batch 3 — Gallica views 41–43 (the last three views of the volume).
% Pass: Opus 5.5 (claude-opus-5-5), 2026-09-22 — first pass, unchecked against the views by a human.
%
% What the three views hold. All three carry writing; none is binding only.
%
%   v. 41   The verso of the leaf whose recto is view 40 (batch 2): an address
%           panel, written upside down with respect to the scan, « Pour le
%           Reuerend Pere Mersenne ». The leaf is folded in four; the rest of
%           the page is blank but for faint show-through. The letter this
%           address closed is on the views before (batch 2); view 40, looked at
%           only as a thumbnail, is a blank recto with a figure at the top right
%           and an ink blot.
%   v. 42   A complete short letter on one page, unsigned and undated, opening
%           « Mon R. Pere ». The writer asks the correspondent to tell « le
%           P. Fabri » that the fifth proposition of his « Theses de geometrie »
%           can be stated much more generally (a segment of a circle equals a
%           sector plus a smaller segment, the two arcs being equal), that the
%           seventh is false (a plane through a given point of a cone that cuts a
%           parabola can be varied to cut infinitely many), and to put to him a
%           proposition on the foci of a hyperbola cut from a scalene cone —
%           which, the writer says, « le P. Grinbergerus » demonstrated — or, if
%           something harder is wanted, « ce beau problesme de M\textsuperscript{r} Roberval »
%           on the surface of an oblique cylinder. Two pen figures in the left
%           margin: a circle lettered g m i h k l, and a triangle through the
%           axis of a cone lettered A B D i c E f h g. The letter ends with two
%           long pen strokes and no signature. Nothing in the pass attributes it.
%   v. 43   The verso of that leaf: a three-line archival note at the top left,
%           « Liasse de 12 cahy / cotées depuis … jusques / a … », and the word
%           « fin » near the foot, both in a hand other than the letter's. The
%           rest is show-through of view 42, mirrored (checked: the slant runs
%           backwards and the circle of the first figure shows at the right).
%           This is the last view of the volume; no back board is filmed.
%
% The hand of v. 42. A rapid, sloping French cursive of the seventeenth
% century, legible at band scale; long s is transcribed as s, said once here.
% The writer closes paragraphs with a point and an oblique stroke (« ./ »),
% which is not reproduced. Capitals are erratic and kept where clear (« Voicy »
% inside the sentence). Spelling of the time kept unflagged: pouuez, esgale,
% circonferance, mesme, fausse, mesner/mesné, patiance, escrire, enuoy-,
% superficie, problesme, « un hyperbole », « une cone ». « Mon R. Pere »,
% « P. », « propos. », « M\textsuperscript{r} » are kept as written. The figures
% and the letters of the figures are set in math italic.
%
% Numbers on the leaves. View 42 carries two figures in the top right corner:
% « 54 » at the very top, small and thin, and below it « 20 », larger and
% heavier. Both are in the position of a foliation and the microfilm, in black
% and white, cannot tell pencil from ink. With 43 views to the volume, « 20 »
% would fit a count of this volume's own leaves and « 54 » an older or wider
% series; but that is an inference, and three views cannot establish either
% series, so no \folio{} is written. View 40 (batch 2's) carries a figure read
% at thumbnail scale as « 52 » or « 62 ». Views 41 and 43 carry no number. The
% round stamp « BIBLIOTHEQUE NATIONALE MSS » (legend partly read) is on view 42,
% in the left third of the page.

\documentclass[11pt,a4paper]{article}
% The preamble shared by every transcription of the mathematicians' manuscripts in Gallica.
%
% It rests on one idea: **uncertainty compounds, it does not resolve.** A
% transcription that smooths over a doubtful word has destroyed the only thing
% it was made for — a reader must be able to tell, at every point, what was on
% the page and what was supplied. The apparatus macros below are therefore the
% critical apparatus, not decoration.
%
% The same commands are recognised by `scripts/render.mjs`, which produces the
% site's reading view, and by `scripts/tei.mjs`. A macro added here must be
% added there too, or rendering fails loudly — which is the wanted behaviour.
%
% Comments here are English, like the rest of the repository. The documents
% this preamble sets are French, and that is deliberate: see the skills.

\ProvidesPackage{germain}[2026/09/15 Transcriptions of mathematicians' manuscripts in Gallica]

\usepackage{amsmath,amssymb,amsthm}
\usepackage{mathrsfs}
% Kept from the parent preamble so the renderer's diagram support stays
% available; these manuscripts draw no commutative diagrams, and nothing here uses it.
\usepackage{tikz-cd}
\usepackage[english,french]{babel}
\usepackage{fontspec}

% --- Greek ------------------------------------------------------------------
% The text face stays fontspec's default, Latin Modern, which has no Greek:
% XeTeX drops every glyph it lacks with a line in the log and nothing on the
% page, and the Greek words the manuscripts quote vanished from the PDFs. So
% the Greek and Greek Extended blocks get a XeTeX character class of their own,
% and entering or leaving a run of it switches the family to CMU Serif — the
% same Computer Modern design, with polytonic Greek — and back. Only the
% family changes: a Greek word in \textit or \textbf keeps its shape and series.
% The transitions to and from every other class carry the tokens that class
% has towards an ordinary letter, so French spacing before « ; » or inside
% « » still holds after a Greek word. No grouping, so no brace or \endgroup
% can come out unbalanced; maths is left alone, since XeTeX inserts nothing
% there. Kept identical in `exercices.sty`.
\makeatletter
\newfontfamily\germainGreekfont{cmun}[
  Extension=.otf, NFSSFamily=germaingreek,
  UprightFont=*rm, ItalicFont=*ti, SlantedFont=*sl,
  BoldFont=*bx, BoldItalicFont=*bi, BoldSlantedFont=*bl]
\newXeTeXintercharclass\germain@greek
\def\germain@greekrange#1#2{%
  \@tempcnta=#1\relax
  \loop
    \XeTeXcharclass\@tempcnta=\germain@greek
    \ifnum\@tempcnta<#2\advance\@tempcnta\@ne\repeat}
\germain@greekrange{"0370}{"03FF}
\germain@greekrange{"1F00}{"1FFF}
\def\germain@greekprev{\rmdefault}
\def\germain@greekin{%
  \edef\germain@greekprev{\f@family}\fontfamily{germaingreek}\selectfont}
\def\germain@greekout{\fontfamily{\germain@greekprev}\selectfont}
\def\germain@greekjoin{%
  \XeTeXinterchartokenstate=\@ne
  \@tempcnta=\z@
  \loop
    \ifnum\@tempcnta=\germain@greek\else
      \XeTeXinterchartoks\germain@greek\@tempcnta=\expandafter{%
        \expandafter\germain@greekout\the\XeTeXinterchartoks\z@\@tempcnta}%
      \XeTeXinterchartoks\@tempcnta\germain@greek=\expandafter{%
        \the\XeTeXinterchartoks\@tempcnta\z@\germain@greekin}%
    \fi
    \ifnum\@tempcnta<4095\advance\@tempcnta\@ne\repeat}
% babel-french sets its own classes, and saves and restores the interchar
% state, each time a language is selected: join again after it does.
\addto\extrasfrench{\germain@greekjoin}
\addto\extrasenglish{\germain@greekjoin}
\AtBeginDocument{\germain@greekjoin}
\makeatother

\usepackage{geometry}
\geometry{a4paper,margin=25mm,marginparwidth=28mm}
\usepackage{marginnote}
\usepackage{xcolor}
\usepackage[normalem]{ulem}
\setlength{\emergencystretch}{3em}
\usepackage{hyperref}
\hypersetup{colorlinks=true,linkcolor=black,urlcolor=black,pdfborder={0 0 0}}

% --- Batch metadata ---------------------------------------------------------
% Taken as they stand by the rendering script, which shows them at the head of
% the reading view. They fix what the file claims to cover. The macro names are
% the parent project's, so that the scripts stay shared; \folder here means the
% volume's slug — `fr-9115` — and \pages counts Gallica views.

\newcommand{\folder}[1]{\gdef\grothFolder{#1}}
\newcommand{\batch}[1]{\gdef\grothBatch{#1}}
\newcommand{\pages}[2]{\gdef\grothFirst{#1}\gdef\grothLast{#2}}
\newcommand{\foldertitle}[1]{\gdef\grothFoldertitle{#1}}

% The holder's own shelfmark, printed as they write it: « Français 9115 ».
\newcommand{\shelfmark}[1]{\gdef\grothShelfmark{#1}}

% The dating, copied verbatim from the holder's catalogue — never inferred
% here. « XIXe siècle » is the BnF's; a pass that dates a leaf from its content
% says so in a \note{}, as its own inference.
\newcommand{\dating}[1]{\gdef\grothDating{#1}}

% The status notice, printed in the title block and repeated as a diagonal
% watermark on every page of the PDF. The manuscripts are in the public
% domain; what this line declares is not a right but a quality — an
% unverified first pass by a machine. The skills write the line; a file
% without it gets no watermark, so leaving it out is visible in the output.
\newcommand{\watermark}[1]{\gdef\grothWatermark{#1}}

\gdef\grothFolder{?}\gdef\grothBatch{}\gdef\grothFirst{?}\gdef\grothLast{?}%
\gdef\grothFoldertitle{}\gdef\grothShelfmark{}\gdef\grothDating{}\gdef\grothWatermark{}

% --- The résumé -------------------------------------------------------------
\newenvironment{resume}
  {\small\setlength{\parskip}{0.45em}}
  {\par\medskip\hrule\medskip}

% --- Keywords ---------------------------------------------------------------
% In English, closing the modernised reading's résumé: the modern vocabulary
% under which to search for what the leaves construct. The manifest extracts
% them to tag the volume — this line is the single source of the tags.
\newcommand{\keywords}[1]{\par\smallskip\noindent
  {\footnotesize\textsc{Keywords} --- \textit{#1}}\par}

% --- The critical apparatus -------------------------------------------------

% The Gallica view number, in the margin. This is the anchor that lets a
% disagreement over a formula be settled by looking at *one* image rather than
% a volume of 750. The site uses it to turn the facsimile as the transcription
% is scrolled. A view is one scanned image: usually one side of a leaf.
\newcommand{\page}[1]{%
  \marginnote{\footnotesize\color{gray!70}\textsf{v.\,#1}}%
  \label{page:#1}\ignorespaces}

% The BnF's pencilled foliation, where the leaf carries it and a pass can read
% it: « 198r ». This is what the literature cites — « ff. 198r–208v » — and
% the one line that turns a citation into a view. Recorded, never inferred:
% a folio a pass cannot read is not written.
\newcommand{\folio}[1]{%
  \marginnote{\footnotesize\color{gray!50}\textsf{f.\,#1}}\ignorespaces}

% The modernised reading's coarser counterpart to \page{}: which views a
% section is drawn from.
\newcommand{\pagerange}[2]{%
  \marginnote{\footnotesize\color{gray!70}\textsf{v.\,#1--#2}}%
  \ignorespaces}

% Illegible: never guessed.
\newcommand{\ill}{\textcolor{red!60!black}{[\,\ldots\,]}}

% A doubtful reading: the word is given, and flagged as uncertain.
\newcommand{\uncertain}[1]{\uline{#1}}

% An editorial addition — a missing letter, an implied word.
\newcommand{\add}[1]{\textcolor{blue!50!black}{[#1]}}

% Struck out by the writer. What was crossed out is part of the document.
\newcommand{\struck}[1]{\sout{#1}}

% The transcriber's note, on the state of the page or the mathematics.
\newcommand{\note}[1]{\par\smallskip{\small\color{gray!60!black}\textit{[#1]}}\par\smallskip}

% A marginal note of hers, distinct from the body of the text.
\newcommand{\marginal}[1]{\par\smallskip\noindent\hspace{1em}%
  {\small\color{gray!60!black}#1}\par\smallskip}

% --- Title ------------------------------------------------------------------
% The title block is French, like the documents themselves.

\AddToHook{shipout/background}{%
  \ifx\grothWatermark\empty\else
    \begin{tikzpicture}[remember picture, overlay]
      \node[rotate=52, scale=3.4, align=center, text=gray!18,
            font=\sffamily\bfseries]
        at (current page.center) {\grothWatermark};
    \end{tikzpicture}%
  \fi}

\AtBeginDocument{%
  \begin{center}
    {\large\bfseries Manuscrits de mathématiciens (Gallica) ---
      \ifx\grothShelfmark\empty\grothFolder\else\grothShelfmark\fi}\\[2pt]
    {\itshape\grothFoldertitle}\\[4pt]
    {\small vues \grothFirst--\grothLast
      \ifx\grothBatch\empty\else\ (lot \grothBatch)\fi}\\[2pt]
    \ifx\grothDating\empty\else
      {\small\color{gray!60!black}Datation du catalogue : \grothDating}\\[2pt]
    \fi
    \ifx\grothWatermark\empty\else
      {\small\bfseries\color{red!45!gray}\grothWatermark}\\[2pt]
    \fi
    {\footnotesize\color{gray!60!black}Transcription établie d'après les images IIIF de Gallica.
     Source gallica.bnf.fr / Bibliothèque nationale de France.\\
     Les lectures incertaines sont soulignées, les illisibles notées [\,\ldots\,],
     les ajouts entre crochets ; \textsf{v.} numérote les vues de Gallica,
     \textsf{f.} la foliotation au crayon de la BnF.}
  \end{center}
  \vspace{1em}}

\endinput


\folder{naf-5161}
\shelfmark{NAF 5161}
\batch{3}
\pages{41}{43}
\dating{1601-1700}
\watermark{Lecture automatique\\non vérifiée}
\foldertitle{Lettres originales de DESCARTES adressées au Père Mersenne (1638-1647), et mémoires divers relatifs aux doctrines de Descartes. II Mémoires divers relatifs aux doctrines de Descartes.}

\begin{document}

\page{41}
\note{Verso du feuillet dont la vue 40 est le recto : le panneau d'adresse de la lettre qui précède, écrit tête-bêche par rapport à la prise de vue, au milieu du feuillet plié en quatre. Le reste de la page est blanc ; des lignes d'écriture y transparaissent sans qu'aucune soit tracée sur cette face. Aucun numéro.}

Pour le Reuerend Pere Mersenne
\note{sur une ligne, le P initial en très grande boucle, un trait de plume au-dessus de « Reuerend ».}

\page{42}
\note{Une lettre entière sur une page, sans date ni signature. En haut à droite, deux chiffres : « 54 », petit et fin, au bord supérieur, et au-dessous « 20 », plus grand et plus appuyé ; aucun n'est retenu comme foliation (voir l'en-tête). Le cachet rond « BIBLIOTHEQUE NATIONALE MSS » est apposé dans le tiers gauche, sur les lignes « que le P. Grinbergerus » à « Il faut monstrer ». Deux figures à la plume dans la marge de gauche.}

Mon R. Pere

vous pouuez mander au P. Fabri. Premierement que la 5.\textsuperscript{e} propos. des Theses de geometrie peut estre enoncee bien plus generalement, Comme par exemple soit quelconque section $gmh$, cette section est esgale au secteur $gli$, et a la petite section $hk$, la \ill{} circonferance $hk$, estant posée esgale, a la \ill{} circonferance $hi$.
\note{deux mots interlignés, écrits petit et serrés, n'ont pas été lus. Le premier (sept ou huit lettres, le haut barré par la traverse du t de « est » de la ligne précédente ; on croit voir « \uncertain{grandes} », que le sens n'appuie pas) est placé au-dessus de « la circonferance $hk$ » ; le second, d'une huitaine de lettres terminées par un long trait, au-dessus de « circonferance $hi$ ». Leur point d'insertion exact n'est pas marqué.}
\note{Figure dans la marge, en regard : un cercle de centre $l$, le diamètre $gk$ horizontal ; sur l'arc supérieur, $m$ à gauche, puis $i$, $h$, $k$ ; tracés la corde $gh$, le rayon $li$, et les cordes $ih$ et $hk$. L'aire du segment $gmh$ y vaut celle du secteur $gli$ augmentée du petit segment $hk$ lorsque l'arc $hk$ égale l'arc $hi$ : c'est l'énoncé de la lettre.}

Et secondement que la 7.\textsuperscript{e} de la mesme geometrie est fausse, car si d'un point donné dans le cone l'on peut mener un plan qui face une parabole, l'on en pourra mener qui en feront un infinité d'autres.
\note{« Et secondement » est souligné d'un long trait qui prolonge le t final ; ce n'est pas une rature.}

\uncertain{Pour ce} que je vous priois qu'il vous demonstrast \ill{} cela n'est pas bien difficile, et vous ne luy demanderez pas s'il vous plaist comme tel, car ie croirois, si ie me voulois donner la patiance que i'en viendrois bien a bout, mais vous luy direz que vous croyez que le P. Grinbergerus l'a demonstré, et que n'ayant pas ce livre, vous le priez de vous en escrire la demonstration, ou bien de vous enuoy\textsuperscript{er} la sienne; Elle est telle.
\note{« ce » est écrit d'une seule lettre qui ressemble à un a ; le P initial est souligné. Après « demonstrast », en fin de ligne, un signe bref — une haste descendante, un i pointé et un trait — n'a pas été lu : ponctuation ou mot court. La fin « er » d'« enuoyer » est écrite au-dessus de la ligne. Dans « Elle est telle », un point et une boucle précèdent « telle ».}

Il faut \struck{\ill{}} monstrer que si dans \struck{le triangle} une cone scalene $bae$, (car il n'y a pas de difficulté en un cone droit) l'on mesne un plan qui fasse un hyperbole de laquelle l'axe soit la ligne \uncertain{$hghfi$}, et le centre $g$, ayant mesné du sommet $f$, une perpendiculaire $fE$, sur l'axe du cone $AED$, et que du centre $g$, et de l'intervalle $gE$, l'on descrive un cercle, ce cercle donnera les foyers dans les points ou il coupera l'axe de l'hyperbole.
\note{entre « faut » et « monstrer », un mot biffé d'un long trait est pris sous le cachet de la bibliothèque et n'a pas été lu. « le triangle », biffé, est écrit au-dessus de la ligne précédente, un signe d'insertion au-dessous ; il est donné ici au point de la ligne où descend ce signe, sous « dans ». Les lettres de l'axe sont nettes une à une sauf la troisième, lue h mais peut-être k : la figure, elle, porte sur cette droite $h$, $g$, $f$ et $i$. Le cône est dit « $bae$ » dans le texte, en minuscules, et « $AED$ » pour l'axe ; la figure porte les majuscules $A$, $B$, $D$, $E$.}
\note{Figure dans la marge, en regard : le triangle par l'axe du cône, de sommet $A$ et de base $BC$ (lettrée $B$, $D$, $i$, $c$ de gauche à droite) ; l'axe $AED$ ; une droite qui coupe le côté $AC$ en $f$ et la base en $i$, prolongée vers le haut jusqu'à $h$, où elle rejoint le côté $BA$ prolongé ; de $f$, la perpendiculaire $fE$ sur l'axe. Près de $A$, une ligne pointillée aboutit à la lettre $g$, précédée d'une marque qui ressemble à un T.}

que si vous desirez luy proposer quelque chose de plus difficile Voicy ce beau problesme de M\textsuperscript{r}. Roberval.

Estant donné un cylindre oblique, retrancher d'un seul trait de compas, sur un cylindre droit, une superficie esgale a la superficie du cylindre oblique donné.
\note{la lettre s'arrête ici, suivie de deux longs traits de plume en guise de clôture, sans formule finale ni signature. Le bas de la page est blanc.}

\page{43}
\note{Verso du feuillet précédent. La lettre de la vue 42 y transparaît en entier, à l'envers. En haut à gauche, trois lignes d'une autre main, à l'encre, une note d'archiviste ; « fin » est écrit seul vers le bas de la page. Aucun numéro.}

Liasse de 12 cahy cotées depuis \uncertain{Y} jusques a \uncertain{L}
\note{sur trois lignes : « Liasse de 12 cahy » / « cotées depuis Y jusques » / « a L ». « cahy » est abrégé, sans doute pour « cahyers ». Les deux cotes sont des lettres isolées, la première un y ou un $\gamma$, la seconde un L ou un C bouclé suivi d'un point ; ni l'une ni l'autre n'est sûre. La note compte les pièces d'un dossier ; elle ne dit pas lesquelles.}

fin
\note{seul, dans le tiers inférieur de la page, au-dessus du pli.}

\end{document}
